\section{The Formalisation}\label{sec:formal}

All results of this article are formalised in Lean~4 \parencite{moura:ullrich:cade-21} with
Mathlib \parencite{mathlib:cpp-20}. The development builds on existing Lean libraries for
deterministic HTN planning and its complexity, from which it takes task networks, decomposition,
progression, the structural classes, encodings of problem descriptions, and the complexity
classes. These libraries also contain the deterministic results that we inherit, such as the
bounds of \textcite{erol:et:al:amai-96} and \textcite{alford:et:al:icaps-15}. Computability
notions come from Mathlib's theory of partial recursive functions \parencite{carneiro:itp-19}.
The FOND development itself comprises about 62{,}500 lines of Lean in 215 files, with about
3{,}500 lemmas and theorems. Every theorem depends only on the three standard axioms of Lean's
classical logic (propositional extensionality, choice and soundness of quotients). There are no
unproved assumptions.

\paragraph{What is stated.}
Each entry of Table~\ref{tab:results} is a Lean theorem about a set of bit strings, namely the
encodings of the problem descriptions of the class that have a solution of the given kind.
Membership is stated as membership in the complexity class, hardness as the existence of a
polynomial-time reduction from every problem in the class, and undecidability, (co-)recursive
enumerability and membership in $\Sigma^1_1$ in terms of Mathlib's computable functions and
recursively enumerable predicates. For $\Sigma^1_1$ we use Kleene's normal form as the definition.
Solution concepts are defined on the semantic level (policies, execution structures, belief
states) exactly as in the papers, with the corrections discussed in
Section~\ref{sec:solutions}. The game characterisations used by the algorithms are proved
equivalent to them. Both readings of strong cyclic solutions are defined, and the literal one is
shown to differ from plan existence on a deterministic example.

\paragraph{What is not formalised.}
Harel's theorem, and the step from it to a fixed two-stack machine in the proof of
Theorem~\ref{thm:sigma11}(iv), are used as cited results; this concerns only
$\Sigma^1_1$-hardness. Neither the analytical hierarchy nor the completeness of ill-founded trees
for $\Sigma^1_1$, on which Harel's proof rests, is currently available in Lean.

\paragraph{Lessons.}
The formalisation was not merely a check of finished proofs. Many of the constructions in this
article were found because a formal proof attempt failed. Examples are the gated cheat tasks of
Theorem~\ref{thm:reg-undec}, the two progress modes of the fair summary game for total order, and
the round counter needed for the fair variant of Theorem~\ref{thm:exptime-hard}. The errors we found
in the literature fall into a few recurring patterns:
\begin{itemize}
  \item gadgets that give the planner an unintended escape (an illegal-move detector that can be
    triggered on purpose, a stop method without an acceptance check);
  \item definitions that are subtly off when read literally (execution sequences that include their
    prefixes, strong cyclic solutions without fairness);
  \item steps that are true but not obvious and were left implicit, such as delaying
    decompositions (Lemma~\ref{lem:delay}), well-definedness of policies on isomorphic nodes,
    preservation of structural classes by compilations, and the treatment of negative
    preconditions.
\end{itemize}
None of the published complexity \emph{claims} that we re-examined turned out to be false, with the
exception of claims about strong cyclic solutions under the literal definition. Several of them,
however, rest on arguments that do not establish them.
